% !TEX program = xelatex
% ============================================================
%  جدول أعمال الاجتماع — Meeting Agenda Template
%  نموذج لإعداد جدول أعمال الاجتماعات بشكل منظم قبل انعقادها
%  Compile with: xelatex main.tex (run twice)
% ============================================================
%  Features:
%    - Header with date, time, location, chair
%    - Time-slot agenda table (Time | Topic | Presenter | Duration)
%    - Objectives section
%    - Pre-meeting reading list
%    - Notes section
% ============================================================

\documentclass[11pt,a4paper]{article}

\usepackage[a4paper,margin=2.5cm,top=2.5cm,bottom=2.5cm,headheight=16pt]{geometry}
\usepackage{graphicx}
\usepackage{array}
\usepackage{booktabs}
\usepackage{tabularx}
\usepackage{multirow}
\usepackage{enumitem}
\usepackage{float}
\usepackage{titlesec}
\usepackage{fancyhdr}
\usepackage{setspace}
\usepackage{xcolor}
\usepackage{amsmath}

\usepackage{bidi}
\usepackage[no-math]{fontspec}
\usepackage[quiet]{polyglossia}

\setmainfont[Scale=1]{NotoKufiArabic-Regular.ttf}
\newfontfamily\arabicfont[Script=Arabic,Scale=0.95]{NotoKufiArabic-Regular.ttf}
\newfontfamily\englishfont[Scale=0.95]{TeX Gyre Termes}

\setdefaultlanguage[calendar=gregorian,hijricorrection=1]{arabic}
\setotherlanguage[variant=british]{english}

\usepackage[breaklinks,hidelinks]{hyperref}

\addto\captionsarabic{%
  \renewcommand{\contentsname}{الفهرس}%
  \renewcommand{\figurename}{الشكل}%
  \renewcommand{\tablename}{الجدول}%
}

\titleformat{\section}{\Large\bfseries}{\thesection}{0.7em}{}
\titlespacing*{\section}{0pt}{1.4ex plus 0.4ex}{0.9ex plus 0.3ex}

\pagestyle{fancy}
\fancyhf{}
\fancyhead[R]{\small\itshape جدول أعمال الاجتماع}
\fancyhead[L]{\small اسم المشروع}
\fancyfoot[C]{\small\thepage}
\renewcommand{\headrulewidth}{0.3pt}

\setlist[itemize]{topsep=0.3em, itemsep=0.15em, parsep=0pt}
\setlist[enumerate]{topsep=0.3em, itemsep=0.15em, parsep=0pt}

\newcolumntype{L}[1]{>{\raggedright\arraybackslash}p{#1}}
\newcolumntype{C}[1]{>{\centering\arraybackslash}p{#1}}

\setstretch{1.3}

\begin{document}

% ============================================================
% TITLE BLOCK
% ============================================================
\begin{center}
{\LARGE\bfseries جدول أعمال الاجتماع}\\[0.3em]
{\large Meeting Agenda}\\[1em]
\end{center}

% ============================================================
% MEETING INFO TABLE
% ============================================================
% TODO: املأ بيانات الاجتماع أدناه.
\begin{table}[H]
\centering
\renewcommand{\arraystretch}{1.4}
\begin{tabularx}{\textwidth}{L{3.5cm} X L{3.5cm} X}
\toprule
\textbf{التاريخ:} & --- & \textbf{الوقت:} & --- \\
\textbf{المكان:} & --- & \textbf{المدة المتوقعة:} & --- \\
\textbf{رئيس الاجتماع:} & --- & \textbf{نوع الاجتماع:} & --- \\
\textbf{اسم المشروع:} & --- & \textbf{رقم الاجتماع:} & --- \\
\bottomrule
\end{tabularx}
\end{table}

\vspace{1em}

% ============================================================
% 1. AGENDA TABLE
% ============================================================
\section{جدول الأعمال}
\label{sec:agenda}

% TODO: عدّل بنود جدول الأعمال والأوقات حسب اجتماعك.
\begin{table}[H]
\centering
\caption{جدول الأعمال بالفترات الزمنية}
\renewcommand{\arraystretch}{1.4}
\begin{tabularx}{\textwidth}{L{2.5cm} X L{3.5cm} C{2cm}}
\toprule
\textbf{الوقت} & \textbf{الموضوع} & \textbf{المقدم} & \textbf{المدة} \\
\midrule
\LR{09:00--09:10} & افتتاح الاجتماع ومراجعة المحضر السابق & رئيس الاجتماع & 10 دقائق \\
\LR{09:10--09:30} & عرض تقدم العمل منذ الاجتماع السابق & مدير المشروع & 20 دقيقة \\
\LR{09:30--10:00} & مناقشة المخاطر والمشكلات الحالية & فريق المشروع & 30 دقيقة \\
\LR{10:00--10:15} & مراجعة الميزانية والمصروفات & المسؤول المالي & 15 دقيقة \\
\LR{10:15--10:30} & استراحة & --- & 15 دقيقة \\
\LR{10:30--10:50} & القرارات المطلوبة والتصويت & رئيس الاجتماع & 20 دقيقة \\
\LR{10:50--11:00} & بنود العمل والخطوات القادمة & مدير المشروع & 10 دقيقة \\
\LR{11:00--11:10} & جدولة الاجتماع القادم والإغلاق & رئيس الاجتماع & 10 دقائق \\
\bottomrule
\end{tabularx}
\end{table}

% ============================================================
% 2. OBJECTIVES
% ============================================================
\section{أهداف الاجتماع}
\label{sec:objectives}

% TODO: اكتب أهداف الاجتماع.
\begin{enumerate}
    \item مراجعة تقدم المشروع والتأكد من سير العمل وفق الخطة.
    \item اتخاذ القرارات اللازمة بشأن المخاطر والمشكلات المطروحة.
    \item الموافقة على الميزانية المحدثة للمشروع.
    \item تحديد بنود العمل وتوزيع المسؤوليات للأسبوع القادم.
\end{enumerate}

% ============================================================
% 3. PRE-MEETING READING
% ============================================================
\section{القراءات قبل الاجتماع}
\label{sec:reading}

% TODO: اذكر المستندات المطلوب قراءتها قبل الاجتماع.
\begin{itemize}
    \item محضر الاجتماع السابق --- مراجعة قبل الاجتماع.
    \item تقرير تقدم العمل الأسبوعي --- مراجعة المؤشرات الرئيسية.
    \item تحديث سجل المخاطر --- مراجعة المخاطر الجديدة.
    \item تقرير الميزانية الحالي --- مراجعة المصروفات والانحرافات.
\end{itemize}

% ============================================================
% 4. NOTES
% ============================================================
\section{ملاحظات}
\label{sec:notes}

% TODO: أضف أي ملاحظات إضافية.
\begin{itemize}
    \item يُرجى الحضور في الموعد المحدد والالتزام بالجدول الزمني.
    \item يُرجى إعداد أي مستندات داعمة قبل الاجتماع وإرسالها للجميع.
    \item يمكن إضافة بنود طارئة لجدول الأعمال بشرط إبلاغ رئيس الاجتماع مسبقاً.
\end{itemize}

\vspace{2em}

\noindent\rule{\textwidth}{0.4pt}

\vspace{0.5em}

\begin{flushleft}
\textbf{أعدّه:} --- \hfill \textbf{التاريخ:} ---
\end{flushleft}

\end{document}
